\documentclass[12pt,a4paper]{article}

\usepackage[utf8]{inputenc}
\usepackage[T2A]{fontenc}
\usepackage[ukrainian]{babel}

\usepackage[
    left=20mm,
    right=20mm,
    top=20mm,
    bottom=20mm
]{geometry}
\usepackage{setspace}
\onehalfspacing
\usepackage{indentfirst}
\usepackage{amsmath}
\usepackage{amssymb}
\usepackage{array}
\usepackage{booktabs}
\usepackage{longtable}
\usepackage{listings}
\usepackage{xcolor}

\lstset{
    language=Python,
    inputencoding=utf8,
    extendedchars=true,
    basicstyle=\ttfamily\small,
    breaklines=true,
    frame=single,
    columns=fullflexible,
    showspaces=false,
    showstringspaces=false,
    keywordstyle=\color{blue},
    stringstyle=\color{green!50!black},
    numbers=left,
    numberstyle=\tiny\color{gray},
    numbersep=8pt
}

\usepackage{graphicx}
\usepackage{float}
\usepackage[
    hidelinks
]{hyperref}
\setlength{\parindent}{1.25cm}
\setlength{\parskip}{0pt}
\setcounter{tocdepth}{1}

\begin{document}


\begin{titlepage}

\begin{center}

КИЇВСЬКИЙ ПОЛІТЕХНІЧНИЙ ІНСТИТУТ
ІМ. ІГОРЯ СІКОРСЬКОГО

\vspace{2cm}

{\large ЗВІТ}\\
{\large ДО ЛАБОРАТОРНОЇ РОБОТИ №2}

\vspace{1cm}

{\large на тему:}

\vspace{0.5cm}

\textbf{<<Реалізація алгоритмів генерації ключів гібридних криптосистем>>}

\end{center}

\vfill

\begin{flushright}
\begin{minipage}{0.4\textwidth}

Виконали:\\
студентки групи ФІ-62 мн\\
Дегтярьова Марія\\
Тимошенко Марія

\end{minipage}
\end{flushright}

\vfill

\begin{center}
Київ -- 2026
\end{center}

\end{titlepage}

\tableofcontents
\newpage

Мета роботи: дослідження принципів формування криптографічно стійких випадкових 
послідовностей, генерації ключового матеріалу, тестування та генерації простих чисел 
для асиметричних криптосистем, а також основних вимог до життєвого циклу 
криптографічних ключів. 

\section{Дослідження принципів генерації криптографічних ключів}

\subsection{Що таке ентропія у контексті генерації криптографічних ключів? Які джерела ентропії можуть використовувати сучасні програмно-апаратні системи?}

Ентропія у контексті криптографії характеризує ступінь невизначеності випадкових даних, які використовуються для створення криптографічних ключів. 
Чим вища ентропія, тим складніше зловмиснику передбачити значення ключа або підібрати його перебором. Навіть якщо ключ має достатню довжину, його безпека може бути низькою, 
якщо він згенерований на основі передбачуваних даних.

Сучасні програмно-апаратні системи можуть використовувати різні джерела ентропії. До них належать фізичні процеси, зокрема тепловий шум електронних компонентів 
і нестабільність частоти коливань електронних схем. Також можуть враховуватися часові характеристики системних подій, рухи миші, натискання клавіш та інші непередбачувані події. 
Операційні системи можуть поєднувати дані з кількох джерел, обробляти їх і використовувати для ініціалізації криптографічно стійких генераторів випадкових чисел. 
При цьому важливо не лише збирати випадкові дані, а й оцінювати їхню якість та рівень ентропії \cite{RFC4086,NIST80090B}.

\subsection{У чому полягає різниця між TRNG, PRNG та криптографічно стійким генератором псевдовипадкових чисел (CSPRNG)?}

TRNG (True Random Number Generator) є генератором справжніх випадкових чисел, який використовує фізичні процеси для отримання випадкових значень. 
Такі генератори можуть спиратися на тепловий шум або інші фізичні явища, результати яких важко передбачити. Їхньою особливістю є те, що випадковість формується 
на основі фізичного джерела, а не лише обчислень.

PRNG (Pseudorandom Number Generator) є генератором псевдовипадкових чисел, який за певним алгоритмом формує послідовність значень на основі початкового стану. 
Такі послідовності можуть мати хороші статистичні властивості, однак за однакових початкових умов генератор відтворює ті самі значення.

CSPRNG (Cryptographically Secure Pseudorandom Number Generator) є криптографічно стійким генератором псевдовипадкових чисел. Він також використовує алгоритм 
для формування послідовності, але має додаткову вимогу щодо непередбачуваності результатів. Знання попередніх значень не повинно давати змоги ефективно визначити 
наступні значення або відновити внутрішній стан генератора за умови дотримання вимог до його реалізації та ініціалізації. 
CSPRNG зазвичай використовує достатньо випадкове початкове значення, отримане з надійного джерела ентропії. Саме такі генератори застосовують для створення криптографічних 
ключів та інших секретних даних \cite{RFC4086,NIST80090A}.

\subsection{Яку роль виконує початкове значення (seed) генератора псевдовипадкових чисел? Які наслідки може мати можливість передбачити або відтворити seed?}

Початкове значення (seed) використовується для ініціалізації генератора псевдовипадкових чисел. На його основі формується початковий стан, 
з якого алгоритм надалі генерує послідовність значень. Для звичайного детермінованого генератора однакове початкове значення за однакових умов призводить 
до відтворення тієї самої послідовності. Це є корисним під час моделювання та проведення експериментів, оскільки дозволяє повторити отримані результати.

У криптографії можливість передбачити або відтворити seed становить серйозну загрозу безпеці. Якщо зловмисник знає початкове значення або може перебрати невелику 
кількість можливих варіантів, він може відтворити послідовність генератора та визначити створені на її основі ключі. У такому випадку навіть ключ великої довжини може виявитися вразливим, 
оскільки фактична кількість можливих ключів обмежується невизначеністю початкового значення. Тому для криптографічних генераторів важливо використовувати достатньо випадковий seed, 
захищати внутрішній стан і правильно реалізовувати механізми ініціалізації та оновлення стану \cite{RFC4086,NIST80090A}.

\subsection{Чому звичайні генератори псевдовипадкових чисел, призначені для моделювання, ігор або статистичних задач, не повинні використовуватися для генерації криптографічних ключів?}

Звичайні генератори псевдовипадкових чисел призначені насамперед для отримання послідовностей, які мають необхідні статистичні властивості та підходять для моделювання, 
ігор або статистичних експериментів. Проте такі генератори не обов'язково забезпечують захист від спроб передбачити наступні значення або відновити їх за відомими результатами. 
Тому статистично випадкова на вигляд послідовність не обов'язково є безпечною для використання в криптографії.

Якщо криптографічний ключ створено за допомогою звичайного генератора, зловмисник може скористатися особливостями алгоритму, передбачуваним початковим значенням 
або відомими результатами його роботи, щоб звузити множину можливих ключів. Це може призвести до розкриття зашифрованих даних або компрометації системи автентифікації. 
Наприклад, у мові Python модуль \texttt{random} призначений для моделювання та статистичних задач, тоді як модуль \texttt{secrets} надає засоби для отримання криптографічно 
стійких випадкових значень. Отже, для генерації криптографічних ключів потрібно використовувати спеціально призначені криптографічно стійкі генератори, які спираються 
на належне джерело випадковості та відповідають вимогам безпеки \cite{RFC4086,PythonRandom,PythonSecrets}.


\subsection{Яким чином сучасні операційні системи надають програмам криптографічно стійкі випадкові значення? Навести приклад відповідного API для однієї операційної системи або мови програмування.}

\subsection{Чим принципово відрізняється генерація випадкового симетричного ключа від генерації пари ключів асиметричної криптосистеми?}

\subsection{Чи гарантує використання алгоритму AES-256 рівень захисту у 256 біт, якщо ключ отримано з недостатньо випадкового або передбачуваного джерела? Обґрунтувати відповідь.}

\subsection{Які основні етапи життєвого циклу криптографічного ключа необхідно враховувати, крім безпосередньої генерації ключа?}

\section{Дослідження генерації простих чисел}

\subsection{Для чого при генерації ключів RSA необхідні великі прості числа?}
а
 
\subsection{Чому на практиці для великих чисел використовуються ймовірнісні тести простоти?}

\subsection{У чому полягає основна ідея тесту Ферма та чому позитивний результат цього тесту не гарантує простоту числа?}

\subsection{Яку проблему тесту Ферма частково вирішують більш досконалі ймовірнісні тести?}

\subsection{Який принцип роботи тесту Міллера--Рабіна та яке значення має кількість його ітерацій?}

\subsection{Чому дуже мала ненульова ймовірність помилки тесту простоти може бути прийнятною при практичній генерації криптографічних ключів?}

\subsection{Які вимоги, крім достатньої довжини, висуваються до простих чисел, що використовуються при формуванні RSA-ключів?}

!!!коротко порівняти
\section{Практичне дослідження генерації випадкових значень}


\subsection{Використання звичайного генератора псевдовипадкових чисел}

Для дослідження відтворюваності псевдовипадкових чисел було використано стандартний модуль \texttt{random} 
мови програмування Python. Спочатку було встановлено фіксоване початкове значення \texttt{seed = 42}, 
після чого згенеровано послідовність із десяти випадкових цілих чисел у діапазоні від 0 до 255. 
Далі генерацію було повторено з тим самим початковим значенням, а отримані послідовності порівняно між собою.

\begin{lstlisting}[language=Python]
import random

seed = 42
n = 10

random.seed(seed)
sequence1 = [random.randint(0, 255) for _ in range(n)]

random.seed(seed)
sequence2 = [random.randint(0, 255) for _ in range(n)]

print("I try:", sequence1)
print("II try:", sequence2)
print("Same?", sequence1 == sequence2)
\end{lstlisting}

У результаті експерименту було встановлено, що обидві послідовності повністю збігаються. 

\begin{lstlisting}[language=Python]
I try: [57, 12, 140, 125, 114, 71, 52, 44, 216, 16]
II try: [57, 12, 140, 125, 114, 71, 52, 44, 216, 16]
Same? True
\end{lstlisting}

Це пояснюється тим, що генератор \texttt{random} є детермінованим і за однакового початкового 
значення та однакових умов генерує ту саму послідовність чисел.

Якби отримані значення використовувалися для формування криптографічних ключів, 
така відтворюваність могла б становити загрозу безпеці. Якщо зловмисник знає початкове значення 
або може його вгадати, він може повторити роботу генератора та отримати ті самі значення, 
що використовувалися під час створення ключа. Це дозволило б звузити пошук можливих ключів 
і потенційно призвести до розкриття захищених даних. Тому звичайний генератор \texttt{random}, 
призначений для моделювання та статистичних задач, не слід використовувати для генерації криптографічних ключів. 
Для цього необхідно застосовувати криптографічно стійкі генератори випадкових чисел.

\subsection{Використання криптографічно стійкого генератора}

Для практичної демонстрації роботи криптографічно стійкого генератора випадкових чисел (CSPRNG)
використано стандартну бібліотеку \texttt{cryptography} у мові Python~\cite{cryptography, SP800-90A}.

Імпорт необхідних модулів:
\begin{lstlisting}[language=Python]
from cryptography.hazmat.primitives.ciphers.aead import AESGCM
\end{lstlisting}

\subsubsection{Генерація криптографічних ключів}

\begin{lstlisting}[language=Python]
key1 = AESGCM.generate_key(bit_length=256)
key2 = AESGCM.generate_key(bit_length=256)

print("Key 1 (hex):", key1.hex())
print("Key 2 (hex):", key2.hex())
\end{lstlisting}

Результат запуску:

\begin{lstlisting}[language=bash]
Key 1 (hex): a1e668462ba47835e8f5d20e6753997a9651ad0c05fd492fe37e47b671c90583
Key 2 (hex): 75d8e1eb2917f2c811df249e01661114570551273bb1851d63f7947aaa3115f1
\end{lstlisting}

\subsubsection{Повторний запуск програми}

Під час повторного запуску програми (повторного виконання скрипту) отримано такі резултати:

\begin{lstlisting}[language=bash]
Key 1 (hex): d7d04a1ff68c38309478f8221868df67fa601470edbc2950b176611de874ce25
Key 2 (hex): 85fb6c2041ab72c0c2dcaaf5b4fb05e64712585629c7b864a02755fad8d70312
\end{lstlisting}

\subsubsection{Порівняння результатів з експериментом 3.1}

У звичайному некриптографічному генераторі (розглянутому в п. 3.1) за встановленого фіксованого
початкового значення (\textit{seed}) згенеровані послідовності повністю відтворюються при кожному
перезапуску програми. Натомість при використанні CSPRNG ключі щоразу є повністю унікальними та
непередбачуваними, а їх відтворення чи повторення після перезапуску є неможливим.

\subsubsection{Джерело криптографічно стійких випадкових значень}

Програмний інтерфейс отримує випадкові значення безпосередньо з операційної системи, наприклад, через
криптографічні функції \texttt{BCryptGenRandom} у Windows або системний виклик \texttt{getrandom()} у
Linux~\cite{SP800-90A, cryptography}. Ядро ОС динамічно акумулює фоновий апаратний шум (ентропію) від
переривань пристроїв введення, таймерів і системних подій, гарантуючи високу якість та непередбачуваність
ключового матеріалу.

\section{Дослідження генерації асиметричних ключів}

\subsection{Генерація RSA-ключа довжиною 2048 біт}

\subsubsection{Програмна реалізація}

Для генерації RSA-ключа довжиною 2048 біт використано криптографічну бібліотеку \texttt{cryptography}. Час генерації вимірювався за допомогою функції \texttt{time.perf\_counter()}. Експеримент було проведено тричі.

\begin{lstlisting}[language=Python]
from cryptography.hazmat.primitives.asymmetric import rsa
import time

for i in range(1, 4):
    start = time.perf_counter()

    private_key = rsa.generate_private_key(
        public_exponent=65537,
        key_size=2048
    )

    end = time.perf_counter()
    print(f"Try (2048) {i}: {end - start:.4f} c")
\end{lstlisting}

\subsubsection{Результати генерації}

Під час трьох запусків час генерації RSA-ключа довжиною 2048 біт становив 0,0289 с, 0,0950 с і 0,0439 с відповідно. Середній час генерації дорівнював приблизно 0,0559 с.

\subsection{Генерація RSA-ключа довжиною 4096 біт}

\subsubsection{Програмна реалізація}

Для дослідження впливу довжини ключа на час генерації експеримент було повторено для RSA-ключа довжиною 4096 біт.

\begin{lstlisting}[language=Python]
from cryptography.hazmat.primitives.asymmetric import rsa
import time

for i in range(1, 4):
    start = time.perf_counter()

    private_key = rsa.generate_private_key(
        public_exponent=65537,
        key_size=4096
    )

    end = time.perf_counter()
    print(f"Try (4096) {i}: {end - start:.4f} c")
\end{lstlisting}

\subsubsection{Результати генерації}

Час генерації RSA-ключа довжиною 4096 біт становив 1,2870 с, 0,2607 с і 0,3652 с відповідно. Середній час генерації дорівнював приблизно 0,6376 с. Результати демонструють помітні відмінності між окремими запусками, що можуть бути пов'язані з особливостями пошуку простих чисел та навантаженням на систему.

\subsection{Порівняння часу генерації}

У таблиці наведено результати вимірювання часу генерації RSA-ключів різної довжини.

\begin{table}[H]
\centering
\caption{Результати генерації RSA-ключів}
\label{tab:rsa_generation}
\begin{tabular}{|c|c|c|}
\hline
\textbf{Довжина ключа, біт} &
\textbf{Номер запуску} &
\textbf{Час генерації, с} \\
\hline
2048 & 1 & 0,0289 \\
\hline
2048 & 2 & 0,0950 \\
\hline
2048 & 3 & 0,0439 \\
\hline
4096 & 1 & 1,2870 \\
\hline
4096 & 2 & 0,2607 \\
\hline
4096 & 3 & 0,3652 \\
\hline
\end{tabular}
\end{table}

Середній час генерації ключа довжиною 2048 біт становив приблизно 0,0559 с, а ключа довжиною 4096 біт — 0,6376 с. Таким чином, у проведеному експерименті генерація ключа довжиною 4096 біт у середньому тривала приблизно в 11,4 раза довше.

\subsection{Аналіз результатів}

Результати експерименту показують, що збільшення довжини RSA-ключа з 2048 до 4096 біт суттєво впливає на час генерації. Це пов'язано зі збільшенням розміру чисел, з якими виконуються арифметичні операції, та складністю пошуку відповідних простих чисел.

Час генерації може відрізнятися навіть для ключів однакової довжини. Під час кожного запуску використовуються нові випадкові значення, тому кількість спроб, необхідних для знаходження простих чисел, може бути різною. Крім того, на результати вимірювань впливають навантаження на процесор та інші процеси, що виконуються на комп'ютері.

Важливу роль у генерації RSA-ключів відіграє криптографічно стійкий генератор псевдовипадкових чисел (CSPRNG), який забезпечує непередбачуваність значень, що використовуються для формування кандидатів на прості числа. Ці кандидати проходять тести простоти, після чого знайдені прості числа використовуються для побудови ключової пари.

Збільшення довжини RSA-ключа також впливає на швидкість криптографічних операцій. Зокрема, операції розшифрування та створення цифрового підпису зазвичай потребують більше обчислень для довших ключів. Водночас збільшення довжини ключа підвищує стійкість RSA до атак, заснованих на факторизації модуля.

Самостійна реалізація генерації RSA-ключів у прикладному програмному забезпеченні зазвичай недоцільна, оскільки потребує коректної реалізації арифметики великих чисел, надійного джерела випадковості та перевірених методів тестування простоти. Помилки в цих компонентах можуть призвести до вразливості ключів. Тому доцільно використовувати перевірені криптографічні бібліотеки, які реалізують необхідні алгоритми та враховують вимоги безпеки.

\section{Аналіз ситуацій порушення безпеки}


\subsection{Ситуація 1. Передбачувана генерація ключа}

Проблема полягає в тому, що для створення AES-256 ключа використовується поточний час. 
Хоча SHA-256 повертає значення довжиною 256 біт, це не означає, що ключ має 256 біт криптографічної стійкості. 
Якщо зловмисник знає приблизний час створення ключа, він може перебрати можливі значення часу, 
обчислити їхні хеші та знайти потрібний ключ. Таким чином, безпека залежить не від довжини хешу, 
а від кількості можливих початкових значень. Для усунення проблеми необхідно генерувати ключ за допомогою 
криптографічно стійкого генератора випадкових чисел (CSPRNG), який забезпечує достатню непередбачуваність.

\subsection{Ситуація 2. Однаковий seed на різних пристроях}

Якщо однакові пристрої використовують однакове початкове значення для генератора псевдовипадкових чисел, 
вони можуть отримувати однакові послідовності. Якщо ці послідовності використовуються для створення 
криптографічних ключів, різні пристрої можуть мати однакові ключі, що значно знижує безпеку системи. 
У такому випадку компрометація одного пристрою може створити загрозу для інших пристроїв із такими самими ключами. 
Для розв'язання проблеми необхідно використовувати системний CSPRNG, який отримує випадкові дані з надійних джерел операційної системи, 
а не встановлювати однакове фіксоване початкове значення.

\subsection{Ситуація 3. Розголошення правильно згенерованого ключа}

Навіть якщо ключ було створено за допомогою надійного CSPRNG, систему не можна вважати повністю захищеною, 
якщо значення ключа записується до журналу програми. Журнали можуть бути доступними іншим користувачам 
або зберігатися тривалий час, тому зловмисник, який отримає до них доступ, зможе використати ключ для 
розшифрування даних або виконання інших операцій від імені системи. Для усунення проблеми не можна записувати 
секретні ключі до журналів. Якщо ключ уже було розголошено, його потрібно замінити, а доступ до журналів обмежити.


\subsection{Ситуація 4. Недостатньо випадкова генерація RSA-ключа}

Під час створення RSA-ключів джерело випадковості має критично недостатню ентропію, внаслідок чого
різні ключі можуть бути сформовані з пов'язаних або повторно використаних простих чисел. 

Основною проблемою тут є низька якість випадкового матеріалу, який використовується 
для генерації простих множників $p$ та $q$. Якщо у двох незалежних модулів
$N_1 = p_1 \cdot q_1$ та $N_2 = p_2 \cdot q_2$ випадково співпаде хоча б один простий множник
(наприклад, $q_1 = q_2$), зловмисник зможе легко знайти цей спільний дільник за допомогою алгоритму
обчислення найбільшого спільного дільника через формулу $\text{gcd}(N_1, N_2) = q$. Це повністю
компрометує та зламає обидві пари асиметричних ключів, дозволяючи розшифрувати приватний трафік або
підробити цифрові підписи. Способом усунення цієї проблеми є використання виключно апаратно
підкріплених системних криптографічно стійких генераторів (CSPRNG) із високим рівнем ентропії та
попередньою перевіркою згенерованих чисел на унікальність~\cite{SP800-56B, RFC8017}.

\subsection{Ситуація 5. Надмірно тривале використання одного ключа}

Інформаційна система використовує один (\textit{master key}) для захисту даних усіх користувачів протягом
тривалого часу без заміни.

Проблема полягає у відсутності політики ротації ключового матеріалу та порушенні принципів мінімізації
ризиків. Така схема створює загрозу накопичення зашифрованого трафіку для подальшого криптоаналізу та
критично збільшує масштаб наслідків у разі витоку ключа: у разі компрометації єдиного \textit{master key} під
загрозою опиняться абсолютно всі дані користувачів за весь період експлуатації системи. Роль регулярної
ротації ключів полягає у зменшенні обсягу даних, захищених одним ключем, і локалізації наслідків можливої
компрометації, оскільки зловмисник зможе отримати доступ лише до сегмента інформації, зашифрованого
поточним зламаним ключем на конкретному часовому інтервалі~\cite{SP800-57}.

\section{Життєвий цикл криптографічного ключа}

Спрощена схема життєвого циклу криптографічного ключа охоплює такі послідовні етапи: генерація
$\rightarrow$ зберігання $\rightarrow$ використання $\rightarrow$ ротація/заміна $\rightarrow$ відкликання
або знищення. Для забезпечення безпеки кожен із цих етапів вимагає врахування специфічних загроз та
дотримання суворих вимог.

На етапі генерації ключа основною загрозою є використання передбачуваних алгоритмів або недостатньої
ентропії (наприклад, застосування звичайних некриптографічних PRNG чи слабкого початкового значення
\textit{seed}), що дозволяє зловмиснику відтворити ключовий матеріал шляхом повного перебору. Вимогою
безпеки тут є обов'язкове використання апаратно підкріплених криптографічно стійких генераторів
випадкових чисел (CSPRNG) у захищеному середовищі~\cite{SP800-90A}.

На етапі зберігання ключа виникають загрози несанкціонованого доступу до файлової системи,
конфігураційних файлів чи вихідного коду програми, де ключі можуть зберігатися у відкритому вигляді, а
також ризики витоку через дампи пам'яті. Вимогою безпеки є розміщення ключів виключно в ізольованих
шифрованих сховищах з використанням обмеження прав доступу та заборони експорту
(\textit{non-exportable keys})~\cite{SP800-57}.

Під час використання ключа основною загрозою є перехоплення секретних даних в оперативній пам'яті,
атаки за побічними каналами (часові атаки або аналіз кешу процесора), а також випадковий запис ключа до
системних журналів (логів) під час діагностики. Вимогами безпеки на цьому етапі є виконання
криптографічних операцій із застосуванням алгоритмів із постійним часом виконання
(\textit{constant-time implementation}), ретельна фільтрація журналів та оперативне очищення ділянок пам'яті
після завершення обчислень~\cite{SP800-57}.

На етапі ротації або заміни ключа ключовою загрозою є відсутність регулярного оновлення ключового
матеріалу, що призводить до надмірного накопичення зашифрованих даних під захистом одного ключа та
збільшує масштаб можливої компрометації. Вимогою безпеки є впровадження автоматизованої політики
планової ротації за розкладом або після обробки визначеного обсягу інформації з підтримкою
контрольованого перехідного періоду~\cite{SP800-57}.

На завершальному етапі відкликання або знищення основними загрозами є неповне або поверхневе
видалення ключа, внаслідок якого залишки ключової інформації можуть бути відновлені зі стаціонарних
носіїв чи кешу, а також затримка блокування компрометованого ключа системою. Вимогами безпеки тут є
реалізація гарантованого безповоротного затирання пам'яті (процедури \textit{zeroization} або
\textit{explicit\_bzero}) та миттєве внесення скомпрометованих ключів до чорних списків відкликання~\cite{FIPS140-3}.



\section*{Висновки}
\addcontentsline{toc}{section}{Висновки}

Під час виконання лабораторної роботи ми дослідили теоретичні та практичні засади формування
криптографічно стійких випадкових послідовностей, принципи генерації ключового матеріалу та побудови
асиметричних криптосистем. Аналізуючи сутність ентропії та різницю між генераторами псевдовипадкових
чисел, ми з'ясували, що звичайні некриптографічні алгоритми моделювання та статистики є непридатними
для захисту інформації через їхню передбачуваність, тоді як криптографічно стійкі генератори
(\textit{CSPRNG}) через системні API операційної системи забезпечують необхідний рівень
непередбачуваності на основі збору апаратного фонового шуму. 

Дослідження тестування простоти великих чисел показало, що використання високоефективних імовірнісних
алгоритмів, зокрема тесту Міллера--Рабіна з достатньою кількістю ітерацій, дозволяє з мізерною ймовірністю
помилки успішно генерувати прості множники, критичні для асиметричних криптосистем на зразок RSA.
Практична частина роботи наочно продемонструвала відтворюваність звичайних \textit{PRNG} порівняно з
унікальними ключами \textit{CSPRNG}, а також підтвердила, що час генерації RSA-ключів суттєво залежить від
їхньої бітової довжини (наприклад, перехід від 2048 до 4096 біт) та складності обчислювальних процедур. 

Аналіз порушень безпеки виявив, що навіть математично стійкі криптографічні алгоритми стають
вразливими внаслідок використання слабких джерел ентропії, однакових початкових значень (\textit{seed})
на пристроях, витоку секретів у системні журнали або відсутності регулярної ротації ключового матеріалу. З
огляду на це, безпека інформаційної системи вимагає комплексного підходу на всіх етапах життєвого циклу
ключа --- від його апаратно підкріпленої генерації у захищеному середовищі до контрольованого
відкликання та гарантованого знищення (\textit{zeroization}) в оперативній пам'яті.


\section*{Перелік використаних джерел}
\addcontentsline{toc}{section}{Перелік використаних джерел}

% Тут будуть джерела
\begin{thebibliography}{99} 

    \bibitem{RFC4086}
    Eastlake, D., Schiller, J., Crocker, S.
    \textit{Randomness Requirements for Security}.
    RFC 4086, 2005.
    \\
    \url{https://www.rfc-editor.org/rfc/rfc4086}

    \bibitem{NIST80090A}
    Barker, E., Kelsey, J.
    \textit{Recommendation for Random Number Generation Using Deterministic Random Bit Generators}.
    NIST Special Publication 800-90A Revision 1, 2015.
    \\
    \url{https://csrc.nist.gov/pubs/sp/800/90/a/r1/final}

    \bibitem{NIST80090B}
    \textit{Recommendation for the Entropy Sources Used for Random Bit Generation}.
    NIST Special Publication 800-90B, 2018.
    \\
    \url{https://csrc.nist.gov/pubs/sp/800/90/b/final}

    \bibitem{PythonRandom}
    Python Software Foundation.
    \textit{random --- Generate pseudo-random numbers}.
    Python Documentation.
    \\
    \url{https://docs.python.org/3/library/random.html}

    \bibitem{PythonSecrets}
    Python Software Foundation.
    \textit{secrets --- Generate secure random numbers for managing secrets}.
    Python Documentation.
    \\
    \url{https://docs.python.org/3/library/secrets.html}

    \bibitem{RSA}
    Rivest, R. L., Shamir, A., Adleman, L.
    \textit{A Method for Obtaining Digital Signatures and Public-Key Cryptosystems}.
    Communications of the ACM, 1978.
    \\
    \url{https://doi.org/10.1145/359340.359342}

    \bibitem{NIST80056B}
    Barker, E., Chen, L., Roginsky, A., Vassilev, A., Davis, R., Simon, S.
    \textit{Recommendation for Pair-Wise Key-Establishment Schemes Using Integer Factorization Cryptography}.
    NIST Special Publication 800-56B Revision 2, 2019.
    \\
    \url{https://csrc.nist.gov/pubs/sp/800/56/b/r2/final}

    \bibitem{FermatTest}
    \textit{Fermat's Little Theorem}.
    Wolfram MathWorld.
    \\
    \url{https://mathworld.wolfram.com/FermatsLittleTheorem.html}

    \bibitem{MillerRabin}
    Miller, G. L.
    \textit{Riemann's Hypothesis and Tests for Primality}.
    Journal of Computer and System Sciences, 1976.
    \\
    \url{https://doi.org/10.1016/0022-0000(76)90043-8}

    \bibitem{Rabin}
    Rabin, M. O.
    \textit{Probabilistic Algorithm for Testing Primality}.
    Journal of Number Theory, 1980.
    \\
    \url{https://doi.org/10.1016/0022-314X(80)90084-0}

    \bibitem{SP800-90A}
    \textit{NIST SP 800-90A Rev. 1. Recommendation for Random Number Generation Using Deterministic Random Bit Generators}.
    \\
    \url{https://csrc.nist.gov/pubs/sp/800/90/a/r1/final}

    \bibitem{cryptography}
    \textit{cryptography}.
    \\
    \url{https://cryptography.io/en/latest/}

    \bibitem{RFC8017}
    \textit{RFC 8017. PKCS \#1: RSA Cryptography Specifications Version 2.2}.
    \\
    \url{https://www.rfc-editor.org/rfc/rfc8017}

    \bibitem{SP800-57}
    \textit{NIST SP 800-57 Part 1 Rev. 5. Recommendation for Key Management: General Guidance}.
    \\
    \url{https://csrc.nist.gov/pubs/sp/800/57/pt1/r5/final}

    \bibitem{FIPS140-3}
    \textit{FIPS PUB 140-3. Security Requirements for Cryptographic Modules}.
    \\
    \url{https://csrc.nist.gov/pubs/fips/140-3/final}

    \bibitem{HAC}
    \textit{Handbook of Applied Cryptography} / A. Menezes, P. van Oorschot, S. Vanstone. --- CRC Press, 2001 (Chapter 4).
    \\
    \url{https://cacr.uwaterloo.ca/hac/}

    \bibitem{SP800-56B}
    \textit{NIST SP 800-56B Rev. 2. Recommendation for Pair-Wise Key Establishment Schemes Using Integer Factorization Cryptography}.
    \\
    \url{https://csrc.nist.gov/pubs/sp/800/56/b/r2/final}

\end{thebibliography}

\end{document}
